\section*{Bab \ref{ch:g11:quad} --- Fungsi dan Persamaan Kuadrat}

\begin{solution}{exo:g11:quad:1}
$x^2-5x+6$: $\Delta = 25 - 24 = 1$, \omterm{def:g11:quad:discriminant}{akarnya} $\frac{5 \pm 1}{2}$, yaitu $2$
dan $3$.

$2x^2+3x-2$: $\Delta = 9 + 16 = 25$, \omterm{def:g11:quad:discriminant}{akarnya} $\frac{-3 \pm 5}{4}$, yaitu
$\frac12$ dan $-2$.

$x^2+2x+5$: $\Delta = 4 - 20 = -16 < 0$, tanpa penyelesaian real.

$9x^2-6x+1$: $\Delta = 36 - 36 = 0$, satu \omterm{def:g11:quad:discriminant}{akar} (kembar)
$x_0 = \frac{6}{18} = \frac13$.
\end{solution}

\begin{solution}{exo:g11:quad:2}
$f(x) = x^2 - 6x + 11 = (x-3)^2 - 9 + 11 = (x-3)^2 + 2$: titik puncak $(3, 2)$.

$g(x) = -2x^2+4x+1 = -2(x^2 - 2x) + 1 = -2\bigl((x-1)^2 - 1\bigr) + 1
= -2(x-1)^2 + 3$: titik puncak $(1, 3)$.

$h(x) = 3x^2 + 6x = 3(x^2 + 2x) = 3\bigl((x+1)^2 - 1\bigr)
= 3(x+1)^2 - 3$: titik puncak $(-1, -3)$.
\end{solution}

\begin{solution}{exo:g11:quad:3}
$x^2 - 3x - 10$: $\Delta = 9 + 40 = 49$, \omterm{def:g11:quad:discriminant}{akarnya} $\frac{3 \pm 7}{2} = -2$
dan $5$; koefisien utamanya positif, sehingga trinomialnya negatif
tegas di antara kedua \omterm{def:g11:quad:discriminant}{akarnya}: himpunan penyelesaiannya $\intoo{-2}{5}$.

$2x^2 + x + 3$: $\Delta = 1 - 24 = -23 < 0$ dan $a = 2 > 0$: selalu
positif, sehingga himpunan penyelesaiannya $\R$.

$-x^2 + 6x - 9 = -(x-3)^2$: nilainya $\geq 0$ hanya ketika $(x-3)^2 \leq 0$,
yaitu di $x = 3$. Himpunan penyelesaiannya $\{3\}$.
\end{solution}

\begin{solution}{exo:g11:quad:4}
$x^2 - 7x + 6$: $1 - 7 + 6 = 0$, jadi $1$ sebuah \omterm{def:g11:quad:discriminant}{akar}; hasil kali kedua
\omterm{def:g11:quad:discriminant}{akarnya} $6$, sehingga \omterm{def:g11:quad:discriminant}{akar} yang lain adalah $6$.

$2x^2 + 5x - 7$: $2 + 5 - 7 = 0$, jadi $1$ sebuah \omterm{def:g11:quad:discriminant}{akar}; hasil kalinya
$\frac{c}{a} = -\frac72$, sehingga \omterm{def:g11:quad:discriminant}{akar} yang lain adalah $-\frac72$.

$x^2 + 4x - 5$: $1 + 4 - 5 = 0$, jadi $1$ sebuah \omterm{def:g11:quad:discriminant}{akar}; hasil kalinya $-5$,
sehingga \omterm{def:g11:quad:discriminant}{akar} yang lain adalah $-5$.
\end{solution}

\begin{solution}{exo:g11:quad:5}
Kedua bilangan itu memenuhi $x^2 - 14x + 45 = 0$ (\cref{prop:g11:quad:vieta});
$\Delta = 196 - 180 = 16$, \omterm{def:g11:quad:discriminant}{akarnya} $\frac{14 \pm 4}{2} = 9$ dan $5$.

Untuk persegi panjangnya: setengah kelilingnya memberi panjang $+$ lebar
$= 17$, dan luasnya memberi hasil kali keduanya $60$. Keduanya memenuhi
$x^2 - 17x + 60 = 0$; $\Delta = 289 - 240 = 49$, \omterm{def:g11:quad:discriminant}{akarnya}
$\frac{17 \pm 7}{2} = 12$ dan $5$. Jadi persegi panjangnya $12 \times 5$.
\end{solution}

\begin{solution}{exo:g11:quad:6}
$\Delta = (m-2)^2 - 4$. Tepat satu penyelesaian ketika $\Delta = 0$:
$(m-2)^2 = 4$, jadi $m - 2 = \pm 2$, yaitu $m = 0$ atau $m = 4$. Tanpa
penyelesaian ketika $\Delta < 0$: $(m-2)^2 < 4$, yaitu $-2 < m - 2 < 2$,
artinya $0 < m < 4$.
\end{solution}

\begin{solution}{exo:g11:quad:7}
Titik puncaknya berada di $t = -\frac{20}{2 \times (-5)} = 2$, dan
$h(2) = -20 + 40 + 1 = 21$: tinggi \omterm{def:g10:functions:extrema}{maksimumnya} $21$ m, dicapai setelah $2$
detik.

$h(t) \geq 16$ berarti $-5t^2 + 20t - 15 \geq 0$, yaitu (setelah dibagi $-5$,
yang membalik arah pertidaksamaannya) $t^2 - 4t + 3 \leq 0$. \omterm{def:g11:quad:discriminant}{Akar}
$t^2-4t+3$ adalah $1$ dan $3$ ($1 - 4 + 3 = 0$, hasil kalinya $3$), jadi
$t^2 - 4t + 3 \leq 0$ tepat pada $\intcc{1}{3}$: bola itu berada paling
sedikit $16$ m di atas tanah antara waktu $t = 1$ s dan $t = 3$ s.
\end{solution}

\begin{solution}{exo:g11:quad:8}
Dengan $X = x^2$: $X^2 - 13X + 36 = 0$, $\Delta = 169 - 144 = 25$, jadi
$X = \frac{13 \pm 5}{2} = 9$ atau $4$. Kembali ke $x$: $x^2 = 9$ memberi
$x = \pm 3$, dan $x^2 = 4$ memberi $x = \pm 2$. Himpunan penyelesaiannya
$\{-3, -2, 2, 3\}$.
\end{solution}

\begin{solution}{exo:g11:quad:9}
Syaratnya adalah $x^2 < x + 2$, yaitu $x^2 - x - 2 < 0$. Karena
$x^2 - x - 2 = (x+1)(x-2)$ (\omterm{def:g11:quad:discriminant}{akarnya} $-1$ dan $2$), pertidaksamaannya berlaku
tegas di antara kedua \omterm{def:g11:quad:discriminant}{akarnya}: $x \in \intoo{-1}{2}$. Pada sketsanya,
\omterm{def:g11:quad:parabola}{parabola} $y = x^2$ menukik di bawah garis $y = x+2$ tepat di antara kedua
titik potongnya, yaitu $(-1, 1)$ dan $(2, 4)$.
\end{solution}

\begin{solution}{exo:g11:quad:10}
Misalkan $x > 0$ memenuhi $x + \frac1x = \frac{13}{6}$. Setelah dikalikan
$6x > 0$: $6x^2 + 6 = 13x$, yaitu $6x^2 - 13x + 6 = 0$. Lalu
$\Delta = 169 - 144 = 25$ dan $x = \frac{13 \pm 5}{12}$, yang memberi
$x = \frac32$ atau $x = \frac23$. Keduanya positif dan saling berkebalikan,
seperti diharapkan: bilangannya adalah $\frac32$ (atau $\frac23$).
\end{solution}

\begin{solution}{exo:g11:quad:11}
\emph{1.} Titik potongnya berpadanan dengan penyelesaian $x^2 = mx - 1$, yaitu
$x^2 - mx + 1 = 0$, yang \omterm{def:g11:quad:discriminant}{diskriminannya} $\Delta = m^2 - 4$. Dua titik
ketika $\abs{m} > 2$; tepat satu ketika $m = \pm 2$; tidak ada ketika $-2 < m < 2$.

\emph{2.} Untuk $m = 2$: $x^2 - 2x + 1 = (x-1)^2 = 0$, jadi titik sentuh
tunggalnya $(1, 1)$. Untuk $m = -2$: $(x+1)^2 = 0$, titik sentuhnya $(-1, 1)$.
Pada kedua kasus itu garisnya bertemu \omterm{def:g11:quad:parabola}{parabola} di satu titik tanpa
memotongnya --- itulah garis singgung di sana, dan memang \omterm{def:g10:reffunc:affine}{gradiennya} $m = \pm 2$
sama dengan nilai di $x = \pm 1$ dari turunan $2x$ \omterm{def:g10:functions:function}{fungsi} $x^2$
(\cref{ch:g11:deriv}).
\end{solution}

\begin{solution}{pb:g11:quad:1}
\textbf{1.} $x^2 - 5x + 6$: $\Delta = 1$, \omterm{def:g11:quad:discriminant}{akarnya} $2$ dan $3$.
$2x^2 + x - 6$: $\Delta = 1 + 48 = 49$, \omterm{def:g11:quad:discriminant}{akarnya}
$\frac{-1 \pm 7}{4}$: yaitu $\frac32$ dan $-2$. $x^2 + x + 1$:
$\Delta = -3 < 0$: tanpa penyelesaian real.

\textbf{2.} $\Delta = 16 + 48 = 64$: \omterm{def:g11:quad:discriminant}{akarnya} $3$ dan $-1$, sehingga
$2x^2 - 4x - 6 = 2(x - 3)(x + 1)$; \omterm{thm:g11:quad:canonical}{bentuk kanoniknya}
$2(x - 1)^2 - 8$. Bentuk terfaktorkan menampilkan \omterm{def:g11:quad:discriminant}{akarnya}, \omterm{thm:g11:quad:canonical}{bentuk
kanonik} menampilkan titik puncak $(1, -8)$, dan bentuk terjabarkan
menampilkan titik potong sumbu $y$ di $-6$.

\textbf{3.} \omterm{def:g11:quad:discriminant}{Akarnya} $1$ dan $3$, koefisien utamanya negatif:
trinomialnya $\geq 0$ \emph{di antara} kedua \omterm{def:g11:quad:discriminant}{akarnya}, jadi
$x \in \intcc{1}{3}$.

\textbf{4.} Kedua bilangan itu adalah \omterm{def:g11:quad:discriminant}{akar} $x^2 - 7x + 12 = 0$:
yaitu $3$ dan $4$. Jika $2$ sebuah \omterm{def:g11:quad:discriminant}{akar} $x^2 - 5x + c$, \omterm{def:g11:quad:discriminant}{akar} yang lain,
$r$, memenuhi $2 + r = 5$ (jumlahnya), jadi $r = 3$, dan
$c = 2 \times 3 = 6$ (hasil kalinya): Viète membacanya langsung.

\textbf{5.} $\Delta = m^2 - 36$: \omterm{def:g11:quad:discriminant}{akar} kembar untuk $m = \pm 6$
(\omterm{def:g11:quad:discriminant}{akarnya} $\mp 3$); tanpa \omterm{def:g11:quad:discriminant}{akar} untuk $-6 < m < 6$; dua \omterm{def:g11:quad:discriminant}{akar} untuk
$\abs m > 6$.

\textbf{6.} $PF^2 = x^2 + \left(x^2 - \frac14\right)^2
= x^2 + x^4 - \frac{x^2}{2} + \frac{1}{16}
= x^4 + \frac{x^2}{2} + \frac{1}{16}
= \left(x^2 + \frac14\right)^2$. Jadi
$PF = x^2 + \frac14$ (kedua faktornya positif).

\textbf{7.} Jarak dari $P(x, x^2)$ ke garis mendatar
$y = -\frac14$ adalah $x^2 + \frac14$ --- tepat sama dengan $PF$. Jadi
setiap titik pada \omterm{def:g11:quad:parabola}{parabola} itu berjarak sama dari \omterm{pb:g11:quad:1}{fokus} dan
direktriksnya.

\textbf{8.} $PF^2 = x^2 + \left(y - \frac14\right)^2$ dan
$\operatorname{dist}^2 = \left(y + \frac14\right)^2$. Kesamaannya
memberi $x^2 + y^2 - \frac y2 + \frac{1}{16} = y^2 + \frac y2 +
\frac1{16}$, jadi $x^2 = y$: tempat kedudukannya tepat \omterm{def:g11:quad:parabola}{parabola} itu.
Satu titik dan satu garis membangkitkan kurva ini, sebagaimana satu titik
dan satu jarak membangkitkan lingkaran.

\textbf{9.} Untuk $y = \frac{x^2}{4f}$, penjabaran yang sama dengan
$F(0, f)$ memberi $PF^2 = x^2 + \left(\frac{x^2}{4f} - f\right)^2
= \left(\frac{x^2}{4f} + f\right)^2$: jarak ke $F$ sama dengan
jarak ke $y = -f$. Untuk antenanya: $9 = \frac{30^2}{4f}$ memberi
$4f = 100$, $f = 25$ cm: penerimanya tergantung seperempat meter
di atas titik puncaknya --- di luar cekungan setinggi $9$ cm itu, seperti
terlihat pada antena yang sesungguhnya.

\textbf{10.} Muka gelombang yang datang sejajar dengan sumbunya
menempuh jarak yang sama untuk mencapai ketinggian direktriksnya;
kesamaan jarak mengubah ``jarak yang sama ke \omterm{pb:g11:quad:1}{direktriks}'' menjadi
``jarak yang sama ke $F$'': semua isyarat pantulannya tiba di
$F$ \emph{seirama}, saling menguatkan --- di situlah tempat mendengarkan.
Bukti yang jujur memerlukan \omterm{def:g10:reffunc:affine}{gradien} garis singgung di setiap titik ---
yaitu turunan pada \cref{ch:g11:deriv}, yang sudah terintip pada
\cref{exo:g11:quad:11}.

\textbf{11.} $y = x\left(1 - \frac{x}{20}\right)$: jatuh di
$x = 20$ m. Titik puncaknya di tengah jalan: $x = 10$, tingginya
$y(10) = 10 - 5 = 5$ m.

\textbf{12.} $y(4) = 4 - \frac{16}{20} = 3.2$ m: melewati tembok $3$ m
dengan sisa ruang, tetapi menghantam tembok $4$ m.

\textbf{13.} $y(10) = 10 - \frac{100}{40} = 7.5$ m $> 7$: truknya
lewat pada $10$ m dari pusat. $y(15) = 10 - \frac{225}{40}
= 4.375$ m: tidak lewat pada $15$ m.

\textbf{14.} $R(p) = -2p^2 + 120p$: titik puncaknya di
$p = \frac{120}{4} = 30$ euro, dengan pendapatan
$R(30) = 30 \times 60 = 1\,800$ euro.

\textbf{15.} $y(320) = 143 \times \left(\frac{320}{640}\right)^2
= \frac{143}{4} \approx 36$ m di atas titik terendahnya.

\textbf{16.} $x^2 = mx - 1$ memberi $x^2 - mx + 1 = 0$,
$\Delta = m^2 - 4$: dua titik untuk $\abs m > 2$, tidak ada untuk
$\abs m < 2$, dan bersinggungan untuk $m = \pm 2$ dengan titik sentuh
$x = \frac m2 = \pm 1$: garis singgungnya $y = 2x - 1$ di $(1, 1)$ dan
$y = -2x - 1$ di $(-1, 1)$.

\textbf{17.} $x^2 = mx + 1$ memberi $\Delta = m^2 + 4 > 0$ untuk
setiap $m$: selalu dua titik potong. Dari sebuah titik dalam,
setiap garis memotong kurvanya dua kali --- garis singgung hanya dapat
ditarik dari luar, persis seperti pada lingkaran.

\textbf{18.} $y = \frac14$ bertemu $y = x^2$ di $x = \pm\frac12$:
panjang tali busurnya $1$. Untuk $y = \frac{x^2}{4f}$: $y = f$ memberi
$x^2 = 4f^2$, $x = \pm 2f$: panjangnya $4f$ --- lebar fokalnya
empat kali jarak fokusnya (perancang antena menyebutnya
latus rectum).

\textbf{19.} Sebuah titik pada $y = x^2$ adalah $(a, a^2)$; \omterm{def:g10:functions:function}{petanya}
$(X, Y) = (ka, ka^2)$, dan $\frac{X^2}{k} = \frac{k^2a^2}{k} =
ka^2 = Y$: jadi kurva \omterm{def:g10:functions:function}{petanya} tepat $y = \frac{x^2}{k}$.
Memilih $k = 4f$ menjangkau setiap \omterm{def:g11:quad:parabola}{parabola}: semuanya
perbesaran satu kurva induk. (Lingkaran pun semuanya sebangun
--- tetapi elips tidak: kepipihannya berarti, perbesarannya tidak.)

\textbf{20.} Sebuah \omterm{def:g10:algebra:equation}{persamaan} dengan tiga busana yang terbaca; tempat
kedudukan titik yang berjarak sama dari sebuah titik dan sebuah garis;
pemantul yang mengumpulkan sinar sejajar di fokusnya dan kabel yang
memikul lantai jembatan yang merata; lintasan setiap bola yang dilempar;
dan, sampai pada perbesaran, satu kurva tunggal --- satu \omterm{def:g11:quad:parabola}{parabola},
di mana-mana.

\end{solution}
